\subsection{The Gap Addressed by This Thesis}
\label{sec:lr-gap}
% Reframed 20 Sep 2026 (author): the thesis is about the quantum walk as a
% general algorithmic tool. The coin is ONE design dimension examined inside
% that programme, not the contribution. Do not restore the earlier version,
% in which all three observations were about the coin.

Two things stand out from the survey above, and together they define the gap
this thesis addresses.

The first is that the walk's claim to generality rests on results that are
almost never exercised. Childs~\cite{PhysRevLett.102.180501} and Lovett et
al.~\cite{ett2010universal} establish that a walk on a suitably constructed
graph is universal for quantum computation, so that every quantum algorithm
has, in principle, a walk formulation. What follows in the literature is not
a programme built on that statement but a collection of individual
algorithms, each constructed for its own problem and analysed in its own
terms: search on particular families of graphs, element distinctness on the
Johnson graph, a mixer for the quantum approximate optimisation algorithm,
kernels between graphs, a feature map for a variational classifier. Every one
of them is a walk; none is presented as an instance of one primitive, and
almost none reports what the walk cost to build as a circuit or what it
bought over the obvious alternative. Generality is asserted at the level of
a compilation theorem and left untested at the level of practice.

The second is that the walk is not one object but a family with several free
choices, and that where quantitative answers do exist they are unusually
sensitive to choices the literature makes by convention. The lackadaisical
self-loop weight decides whether search on the grid obtains a speed-up at
all, and its correct value shifts with the number of marked
vertices~\cite{Nahimovs_2018}, need not be uniform across
vertices~\cite{Rhodes_2019}, and may break the graph's symmetry without
loss~\cite{Rapoza_2021}. Search on the two-dimensional grid succeeds with
the flip-flop shift and fails outright with the moving
shift~\cite{Ambainis_2004}. The coin itself --- position-dependent,
step-dependent or a general element of $U(d)$ --- is known to change the
walk materially~\cite{Tregenna_2003, Ahmad_2020, Panahiyan_2018, Li_2012,
Souza_2013}, and yet outside the search literature it is fixed in advance to
Hadamard, Grover or Fourier and reported for that choice alone. The
parameters that decide whether a walk-based algorithm works are therefore
the ones least often varied.

This thesis accordingly treats the quantum walk as a general algorithmic
tool and exercises the same primitive across several algorithm families on
common terms: order finding, the quantum approximate optimisation algorithm
and unstructured search in Chapter~\ref{ch:algorithms}, and quantum neural
networks, variational classifiers, quantum kernels and graph classification
in Chapters~\ref{ch:methodology} and~\ref{ch:results}. For each the same
three questions are asked --- what the walk formulation actually is, what it
costs as a circuit, and what it achieves against a tuned baseline --- and
the design choices inside the walk, the coin above all but also the shift
and the formulation, are varied rather than fixed by convention wherever an
application makes the comparison meaningful. The question is not whether one
coin outperforms another on one task; it is what the walk is worth as a way
of building quantum algorithms. All of it is answered on simulators, under
the reporting discipline of Bowles et al.~\cite{Bowles_2024}, and a negative
answer for any family is a result and is reported as one.

